\section{Hypothesis}\label{hypothesis}

\subsection{Formal Statement}\label{formal-statement}

\textbf{H₁ (Privacy):} On the synthetic clinical trial corpus at ε =
3.0, δ = 1e-5, the composed NGSP system (Safe Harbor stripping + neural
router + query synthesis / DP bottleneck) reduces the inversion
attacker's span-recovery F1 to ≤ 0.09 --- the Expert Determination
re-identification threshold established in 45 CFR §164.514(b).

\textbf{H₂ (Utility):} Under the same configuration, the utility ratio
(mean proxy-answer quality / mean original-answer quality as scored by
Claude-as-judge) remains ≥ 0.85.

\textbf{H₃ (Composition):} No single component achieves both bounds
simultaneously. The composed system achieves them through layered
defense: Safe Harbor eliminates structured PHI; query synthesis breaks
the injective link between task intent and entity specifics; DP noise
degrades residual quasi-identifier signal.

\subsection{Operationalisation}\label{operationalisation}

{\def\LTcaptype{none} % do not increment counter
\begin{longtable}[]{@{}
  >{\raggedright\arraybackslash}p{(\linewidth - 6\tabcolsep) * \real{0.2162}}
  >{\raggedright\arraybackslash}p{(\linewidth - 6\tabcolsep) * \real{0.2703}}
  >{\raggedright\arraybackslash}p{(\linewidth - 6\tabcolsep) * \real{0.2973}}
  >{\raggedright\arraybackslash}p{(\linewidth - 6\tabcolsep) * \real{0.2162}}@{}}
\toprule\noalign{}
\begin{minipage}[b]{\linewidth}\raggedright
Symbol
\end{minipage} & \begin{minipage}[b]{\linewidth}\raggedright
Quantity
\end{minipage} & \begin{minipage}[b]{\linewidth}\raggedright
Threshold
\end{minipage} & \begin{minipage}[b]{\linewidth}\raggedright
Source
\end{minipage} \\
\midrule\noalign{}
\endhead
\bottomrule\noalign{}
\endlastfoot
F1\_inv & Inversion attacker token-level F1 (Attack 3) & ≤ 0.09 & Expert
Determination / de-identification guidance, HHS \\
U\_ratio & mean(proxy\_score) / mean(original\_score) & ≥ 0.85 &
Clinical utility floor, project-internal \\
AUC\_mem & Membership inference AUC (Attack 4) & ≤ 0.55 & Random
classifier = 0.5; 0.55 allows ±1σ \\
leak\_lit & Verbatim literal leak rate over sensitive spans (Attack 1) &
≤ 0.01 & Near-zero for structured identifiers post-Safe-Harbor \\
\end{longtable}
}

\subsection{Falsification Conditions}\label{falsification-conditions}

H₁ is falsified if any single attack class produces F1\_inv
\textgreater{} 0.09 on the held-out 20\% evaluation split at ε = 3.0.

H₂ is falsified if the mean utility ratio drops below 0.85 over the
50-document evaluation set.

H₃ is falsified if the ablation shows any single component achieves both
H₁ and H₂ simultaneously --- this would mean composition adds no
marginal value.

\subsection{Scope and Limitations}\label{scope-and-limitations}

This hypothesis is evaluated \textbf{only on synthetic data}. The corpus
contains no real patient records, no real compound codenames, and no
real institutional identifiers. Generalization to real clinical trial
documents is explicitly out of scope and must be assessed in a
separately governed study with appropriate data handling controls.

The threat model is a computationally bounded passive adversary with
access to the proxy text but not to the NGSP internals, entity\_map, or
DP noise seed. Adaptive adversaries with partial system knowledge are
not evaluated.
